\section{The integral parameter comparison}

\subsection{The blind control: local charts and coherent direct images}

The proof in this section was written before opening either version of
arXiv:2610.00250.  Only the construction data of AC26 are used for the
threefold.  Write $B=\mathbb P^1$, let $S_1,S_2$ be the reduced fibres of
multiplicities $3,4$, let $W$ be the cusp fibre, and let $D$ be its reduced
double locus.  Thus $D$ consists of three rational curves joining the two
triple points.  Set
\[
 L=f^*\mathcal O_B(-1)\otimes\mathcal O_X(2S_1+S_2),\quad
 F=\mathcal O_X(S_1+S_2),\quad
 E=\ker(df),\quad J=\mathcal I_D\otimes F.
\]
Let $\widehat B$ denote the analytic root orbifold with orders $3,4$ at
$p_1,p_2$, let $\pi:\widehat B\to B$ be its coarse map, and let
$h:X\to\widehat B$ be defined by the canonical sections of $S_1,S_2$.
In particular $f=\pi h$.  This notation means descent on the ordinary
root discs by the specified finite cyclic actions.  Taking invariants is
exact over $\mathbb C$, by averaging.  Denote by $\widehat L,\widehat F$
the root line bundles whose pullbacks are $L,F$.

\begin{theorem}\label{night:P5.1}\label{stmt:P5.1}
\StatementRecord{night:P5.1}{proved}{Historical events: GPT-6 Astra (OpenAI), ChatGPT Pro mode, effort Pro, S6\_06\_BLIND, 16--18 November 2026: targeted proof check. Also checked in this writing. Pre-continuation writing events: GPT-6 Astra (OpenAI), ChatGPT Pro mode, effort Pro, S6\_10\_PAPERS, 18 November 2026: targeted proof check. Continuation metadata events: GPT-6 Astra (OpenAI), ChatGPT Pro mode, effort Pro, S6\_10\_PAPERS continuation, 19 November 2026: used and rechecked in this run; status-vocabulary reconciliation; metadata reconciliation only; no new mathematical check. \emph{Pre-continuation recorded proof scope.} Read the full cp25 re-derivation and re-derived the local syzygy/root-cohomology, KS period equation, integral translation-sheaf, or order-two obstruction calculation stated in the respective proof. These are credited DGKX comparisons; the event does not assert an independent historical discovery.}{the definitions and exact calculations in this statement}

The sheaf $E$ is locally free of rank two and is the pullback of the
translation Lie bundle $V$ on $\widehat B$.  There is an exact sequence
\begin{equation}\label{eq:P5tangent}
 0\longrightarrow E\longrightarrow T_X\longrightarrow J\longrightarrow0.
\end{equation}
\end{theorem}
\begin{proof}
\begin{enumerate}[label=\arabic*.,leftmargin=*]
\item On a finite root chart $t=s^m$, the image of the differential is
$(s^{m-1})f^*T_B$.  The root chart maps étale to $X$ because its affine
cyclic action is free.  Thus the vanishing orders along $S_1,S_2$ are
$2,3$.

\item At a cusp double point use $q=xy$ in coordinates $(x,y,z)$.  The image of
$dq$ is $(x,y)$, and its kernel is freely generated by
$x\partial_x-y\partial_y$ and $\partial_z$: reducing $ya+xb=0$ modulo
$x$ and $y$ gives $a=xA,b=-yA$.  At a triple point $q=xyz$, reduction of
$yz a+xz b+xy c=0$ modulo $x,y,z$ gives
$a=xA,b=yB,c=zC$, and cancellation gives $A+B+C=0$.  Hence its kernel is
freely generated by $(x,-y,0),(x,0,-z)$ and its image is
$(xy,xz,yz)$, the reduced ideal of the three axes.  At a smooth point of
$W$ the differential is surjective.  Consequently its global image is
\[
 \mathcal I_D\otimes f^*T_B\otimes\mathcal O_X(-2S_1-3S_2)
 =\mathcal I_D\otimes\mathcal O_X(S_1+S_2),
\]
using $3S_1=f^*(p_1)$ and $4S_2=f^*(p_2)$.  This proves the exact sequence
including the scheme structure at a triple point.  On the toric charts
the two vertical logarithmic torus fields give exactly this kernel.
Period translations and collar translations have identity vertical
derivative; finite generators act by their vertical linear matrices.
The evaluations of the translation Lie bundle therefore glue to
$h^*V\simeq E$.
\end{enumerate}
\end{proof}

\begin{theorem}\label{night:P5.2}\label{stmt:P5.2}
\StatementRecord{night:P5.2}{proved}{Historical events: GPT-6 Astra (OpenAI), ChatGPT Pro mode, effort Pro, S6\_06\_BLIND, 16--18 November 2026: targeted proof check. Also checked in this writing. Pre-continuation writing events: GPT-6 Astra (OpenAI), ChatGPT Pro mode, effort Pro, S6\_10\_PAPERS, 18 November 2026: targeted proof check. Continuation metadata events: GPT-6 Astra (OpenAI), ChatGPT Pro mode, effort Pro, S6\_10\_PAPERS continuation, 19 November 2026: used and rechecked in this run; status-vocabulary reconciliation; metadata reconciliation only; no new mathematical check. \emph{Pre-continuation recorded proof scope.} Read the full cp25 re-derivation and re-derived the local syzygy/root-cohomology, KS period equation, integral translation-sheaf, or order-two obstruction calculation stated in the respective proof. These are credited DGKX comparisons; the event does not assert an independent historical discovery.}{P5.1}

On the root orbifold,
\[
 V\simeq\mathcal O\oplus\widehat L,
 \qquad R^qh_*\mathcal O_X\simeq
 \bigwedge^q(\mathcal O\oplus\widehat L)\quad(0\leq q\leq2).
\]
On $B$, the direct images used below are
\[
\begin{array}{c|ccc}
 & f_* & R^1f_* &R^2f_*\\ \hline
 \mathcal O_X&\mathcal O&\mathcal O\oplus\mathcal O(-1)&\mathcal O(-1)\\
 L&\mathcal O(-1)&\mathcal O(-1)^{\oplus2}&\mathcal O(-1)\\
 F&\mathcal O&\mathcal O^{\oplus2}&\mathcal O\\
 E&\mathcal O\oplus\mathcal O(-1)&\mathcal O\oplus\mathcal O(-1)^{\oplus3}
       &\mathcal O(-1)^{\oplus2}\\
 J&\mathcal O(-1)&\mathcal O(-1)^{\oplus2}&\mathcal O .
\end{array}
\]
\end{theorem}
\begin{proof}
\begin{enumerate}[label=\arabic*.,leftmargin=*]
\item Applying the homology matrices of the construction to the last two
period columns gives the vertical linear factors
\[
 R_1=\begin{pmatrix}-1/\tau&0\\(1-\mu)/\tau&1\end{pmatrix},\qquad
 R_2=\begin{pmatrix}1/\tau&0\\-\mu/\tau&1\end{pmatrix}.
\]
The second coordinate line is fixed.  The other eigenvalues at the root
points are $\zeta_1^2,\zeta_2$, where
$\zeta_j=\exp(-2\pi i/m_j)$.  At the cusp both logarithmic torus fields
extend in the abstract bundle $E$ by Statement P5.1.

\item The modular discriminant
$\Delta(\tau)=q_\tau\prod_{n\geq1}(1-q_\tau^n)^{24}$ has weight $12$,
has no zero on the upper half-plane, and has a simple zero at the cusp.
The determinant factor of $R_g$ is the reciprocal modular factor, up to
a sign which disappears in its twelfth power.  Thus $\Delta^{-1}$ is a
meromorphic section of $(\det V)^{12}$ with precisely a simple pole at
$p_0$.  Its two root characters determine
\begin{equation}\label{eq:P5L}
 \det V=\widehat L
 =\pi^*\mathcal O_B(-1)\otimes
   \mathcal O_{\widehat B}(2p_1/3+p_2/4).
\end{equation}
Indeed the right side has the prescribed root characters and its
twelfth power is $\pi^*\mathcal O_B(-1)$; the quotient descends to
$\mathbb P^1$ and has trivial twelfth power, so is trivial.

\item For every root line bundle the invariant direct image is
\begin{equation}\label{eq:P5floor}
 \pi_*\bigl(\pi^*\mathcal O_B(d)\otimes
 \mathcal O_{\widehat B}(ap_1/3+bp_2/4)\bigr)
 =\mathcal O_B\bigl(d+\lfloor a/3\rfloor+\lfloor b/4\rfloor\bigr).
\end{equation}
This follows on a root disc by writing its invariant meromorphic
coefficients in integral powers of $t=s^m$; the largest allowed pole
order is $\lfloor a/m\rfloor$.  It also proves the formula for negative
integers.  The invariant line in $V$ gives
$0\to\mathcal O\to V\to\widehat L\to0$.  Its extension group is
$H^1(\widehat B,\widehat L^{-1})=H^1(B,\mathcal O(-1))=0$, so it splits.

\item On a smooth torus, compute $H^1(\mathcal O)$ as the quotient of flat
cohomology by the two period rows.  In that quotient $\gamma,u$ are a
basis, with
\[
 w=-6\mu\gamma-\tau u,\qquad \delta=-\beta\gamma-\mu u.
\]
The transitions are triangular with diagonal $(1,\det R_g)$; for example
$T_1u=-u-w$ gives second column
$(6\mu(g_1z),\tau(g_1z)-1)^{\mathsf t}$.  The image of $\gamma$ is
nowhere zero: comparison of the $w,\delta$ coefficients in the period
rows shows that $\gamma$ never lies in their span.

\item We check the extension and base change at the cusp.  Its normalization
$S=dP_6$ and its double curves give the exact branch complex
\begin{equation}\label{eq:P5normalization}
0\to\mathcal O_W\to\nu_*\mathcal O_S\to
 \bigoplus_{i=1}^3\mathcal O_{D_i}\to
 \mathbb C_P\oplus\mathbb C_Q\to0.
\end{equation}
At $xy=0$ this is the fibre-product identity; at $xyz=0$ it is
inclusion-exclusion of the three convergent power-series rings.  The
restriction differences have the single alternating compatibility
$g_{12}(0)-g_{13}(0)+g_{23}(0)=0$, which suffices to reconstruct the
three branch functions.  All the strata have vanishing higher coherent
cohomology.  Their global complex is
\[
 \mathbb C\xrightarrow{0}\mathbb C^3
 \xrightarrow{\left(\begin{smallmatrix}1&-1&1\\1&-1&1\end{smallmatrix}\right)}
 \mathbb C^2.
\]
It gives $(h^0,h^1,h^2)(\mathcal O_W)=(1,2,1)$.  It is the cellular
cochain complex of the two-triangle torus.  Its product in degrees
$(1,1)$ is the nonzero exterior product to the oriented top class.

\item The constant branch complex identifies
$H^1(W,\mathbb C)\to H^1(W,\mathcal O_W)$ as an isomorphism.  A loop
recording a component translation in the toroidal cover has its class
in $\Lambda/\Lambda_{\rm tor}$; $\gamma,u$ evaluate as its two
coordinates.  Locally constant Cech cocycles on this cover, viewed as
holomorphic cocycles, therefore extend the two classes through the
cusp, and their central values are the basis just calculated.

\item On each ordinary or root disc, the local ring of the smooth total
space is torsion free over the base DVR, hence flat.  Properness is part
of the construction and every fibre has dimension two.  The structure
sheaf is flat and its fibre cohomology dimensions are the constant
values $(1,2,1)$.  Grauert's cohomology and base-change theorem, \S7
Satz 5, applies.  We checked the 1963 corrigendum, whose correction in
\S7 concerns Satz 2.  Thus the direct images are locally free and commute
with base change.  The extension
$0\to\mathcal O\to R^1h_*\mathcal O_X\to\widehat L\to0$ exists on all
charts and splits by the preceding extension-group calculation.  The
relative cup product identifies $R^2h_*\mathcal O_X$ with $\widehat L$,
being nonzero both on smooth tori and on $W$.  Connected fibres and the
maximum principle give $h_*\mathcal O_X=\mathcal O$.

\item For a line bundle $A$ on $\widehat B$, projection and descent give
\[
 R^qf_*h^*A=\pi_*\bigl(A\otimes
              \bigwedge^q(\mathcal O\oplus\widehat L)\bigr).
\]
On a root disc this means tensoring followed by the averaging
projector, so it includes the nonreduced coarse fibres.  Formula
\eqref{eq:P5floor} gives the first four rows in the statement.

\item For the last row one must compute restriction to $D$.  The conductor
square gives
\[
0\to\mathcal O_W\to\nu_*\mathcal O_S\oplus\mathcal O_D
\to\nu_*\mathcal O_{\partial S}\to0,
\]
where $\partial S$ is the hexagon.  Here $h^1(\mathcal O_{\partial S})=1$
and $h^1(\mathcal O_D)=2$.  The equality conditions at the two triple
points give incidence matrix
\[
 \begin{pmatrix}1&-1&0\\0&1&-1\\1&-1&0\\0&1&-1\end{pmatrix}.
\]
Orient the three theta-graph edges from $P$ to $Q$.  The hexagon
traversal maps to $e_1-e_2+e_3-e_1+e_2-e_3=0$.  The constant-to-coherent
branch comparison is an isomorphism in degree one, so the induced map
$H^1(D,\mathcal O_D)\to H^1(\partial S,\mathcal O_{\partial S})$ is
zero.  The conductor sequence then makes
$H^1(W,\mathcal O_W)\to H^1(D,\mathcal O_D)$ an isomorphism.

\item The canonical section of $F$ is nonzero near $W$, so $F|_D=\mathcal O_D$.
In $0\to J\to F\to\mathcal O_D\to0$, the maps of direct images
$\mathcal O\to\mathbb C_{p_0}$ and
$\mathcal O^{\oplus2}\to\mathbb C_{p_0}^{\oplus2}$ are evaluation and
evaluation followed by this isomorphism.  They are surjective, with
kernels $\mathcal O(-p_0)$ in the indicated ranks.  Finally
$H^2(D,\mathcal O_D)=0$ and $R^2f_*F=\mathcal O$.  This gives the last row.
\end{enumerate}
\end{proof}

\begin{theorem}\label{night:P5.3}\label{stmt:P5.3}
\StatementRecord{night:P5.3}{proved}{Historical events: GPT-6 Astra (OpenAI), ChatGPT Pro mode, effort Pro, S6\_06\_BLIND, 16--18 November 2026: targeted proof check. Also checked in this writing. Pre-continuation writing events: GPT-6 Astra (OpenAI), ChatGPT Pro mode, effort Pro, S6\_10\_PAPERS, 18 November 2026: targeted proof check. Continuation metadata events: GPT-6 Astra (OpenAI), ChatGPT Pro mode, effort Pro, S6\_10\_PAPERS continuation, 19 November 2026: used and rechecked in this run; status-vocabulary reconciliation; metadata reconciliation only; no new mathematical check. \emph{Pre-continuation recorded proof scope.} Read the full cp25 re-derivation and re-derived the local syzygy/root-cohomology, KS period equation, integral translation-sheaf, or order-two obstruction calculation stated in the respective proof. These are credited DGKX comparisons; the event does not assert an independent historical discovery.}{P5.2}

For every admissible $c$,
\[
 \bigl(h^0(T_{X_c}),h^1(T_{X_c}),h^2(T_{X_c}),h^3(T_{X_c})\bigr)
 =(1,1,1,0).
\]
\end{theorem}
\begin{proof}
\begin{enumerate}[label=\arabic*.,leftmargin=*]
\item The Laurent--Cech calculation on the two standard charts of
$\mathbb P^1$ gives $h^0(\mathcal O(d))=\max(d+1,0)$ and
$h^1(\mathcal O(d))=\max(-d-1,0)$.  Leray has only columns zero and one,
so has no higher differential.  Statement P5.2 gives
\[
 H^q(X,E)=(\mathbb C,\mathbb C,0,0),\qquad
 H^q(X,J)=(0,0,\mathbb C,0).
\]
The long exact sequence of \eqref{eq:P5tangent} proves the assertion;
every adjacent group which could change these dimensions is zero.
The bundle and incidence computations are independently recomputed in
\texttt{code/part5/bundle\_cohomology.py}.
\end{enumerate}
\end{proof}

\begin{theorem}\label{night:P5.4}\label{stmt:P5.4}
\StatementRecord{night:P5.4}{proved}{Historical events: GPT-6 Astra (OpenAI), ChatGPT Pro mode, effort Pro, S6\_06\_BLIND, 16--18 November 2026: targeted proof check. Also checked in this writing. Pre-continuation writing events: GPT-6 Astra (OpenAI), ChatGPT Pro mode, effort Pro, S6\_10\_PAPERS, 18 November 2026: targeted proof check. Continuation metadata events: GPT-6 Astra (OpenAI), ChatGPT Pro mode, effort Pro, S6\_10\_PAPERS continuation, 19 November 2026: used and rechecked in this run; status-vocabulary reconciliation; metadata reconciliation only; no new mathematical check. \emph{Pre-continuation recorded proof scope.} Read the full cp25 re-derivation and re-derived the local syzygy/root-cohomology, KS period equation, integral translation-sheaf, or order-two obstruction calculation stated in the respective proof. These are credited DGKX comparisons; the event does not assert an independent historical discovery.}{P5.3}

The germ of the $c$-parameter family at each admissible member is a
miniversal deformation germ.  The Kuranishi germ is a smooth disc.
\end{theorem}
\begin{proof}
\begin{enumerate}[label=\arabic*.,leftmargin=*]
\item The period and gluing formulas depend holomorphically on $c$, and the
strict admissibility condition is open, so a small $c$-disc gives a
proper holomorphic family.  A first-order trivialization of its
deformation in the direction $dc=1$ would carry $f$ to
$f+\epsilon v$, with $v\in H^0(X,f^*T_B)$.  Since $f_*\mathcal O_X$
is $\mathcal O_B$, this is a base vector field.  On a simply connected
regular base chart the integral marking is constant under a first-order
trivialization and is the identity at the reference point.  Translation
does not change periods.  Its infinitesimal linear coordinate change
$H(t)$ must therefore satisfy
\[
 \partial_c\Pi+v(t)\partial_t\Pi=H(t)\Pi.
\]
The last two period columns are the constant identity columns, so
$H=0$.  The upper entry of the second column gives $v\tau'=0$; the
constructed nonconstant $\tau$ has nonzero derivative on the regular
chart, and thus $v=0$ everywhere.  The lower entry of the first column
now says $1=0$.  This proves that the Kodaira--Spencer map is nonzero,
and by P5.3 it is an isomorphism.

\item Douady's \S VII.8 principal theorem gives a finite-dimensional versal
analytic germ, and \S VIII.1 Theorem 1 and its corollary give minimality
and uniqueness of its tangent map.  First-order changes of holomorphic
transition maps are tangent-valued Cech cocycles and changes of
trivialization are coboundaries, so its tangent space here is $H^1(T)$.
Embed the minimal germ as the zero space of an ideal
$I\subset\mathbb C\{z\}$.  Versality gives a map from our disc with
invertible derivative.  Substitution by a convergent series with
nonzero derivative is injective and annihilates $I$, forcing $I=0$.
The inverse-function theorem then identifies the germs.  This also
excludes a nonreduced embedded germ.
\end{enumerate}
\end{proof}

\subsection{All automorphisms and the integral parameter period}

\begin{theorem}\label{night:P5.5}\label{stmt:P5.5}
\StatementRecord{night:P5.5}{proved}{Historical events: GPT-6 Astra (OpenAI), ChatGPT Pro mode, effort Pro, S6\_06\_BLIND, 16--18 November 2026: targeted proof check. Also checked in this writing. Pre-continuation writing events: GPT-6 Astra (OpenAI), ChatGPT Pro mode, effort Pro, S6\_10\_PAPERS, 18 November 2026: targeted proof check. Continuation metadata events: GPT-6 Astra (OpenAI), ChatGPT Pro mode, effort Pro, S6\_10\_PAPERS continuation, 19 November 2026: used and rechecked in this run; status-vocabulary reconciliation; metadata reconciliation only; no new mathematical check. \emph{Pre-continuation recorded proof scope.} Read the full cp25 re-derivation and re-derived the local syzygy/root-cohomology, KS period equation, integral translation-sheaf, or order-two obstruction calculation stated in the respective proof. These are credited DGKX comparisons; the event does not assert an independent historical discovery.}{the definitions and exact calculations in this statement}

For every admissible twist of the $(3,4)$ construction,
\[
 \operatorname{Aut}(X_c)=\mathbb C/\mathbb Z\simeq\mathbb C^*.
\]
The parameter $a\bmod\mathbb Z$ acts by translation by $a\widehat\delta$.
\end{theorem}
\begin{proof}
\begin{enumerate}[label=\arabic*.,leftmargin=*]
\item Taking determinants in \eqref{eq:P5tangent}, whose quotient ideal has
codimension-two support, gives
\[
 K_X^{-1}=L\otimes F
 =f^*\mathcal O(-1)\otimes\mathcal O(3S_1+2S_2)
 =\mathcal O(2S_2),\qquad K_X^{-2}=f^*\mathcal O_B(1).
\]
The complete linear system of $K_X^{-2}$ is exactly $f$, since its
sections are the two sections pulled back from $B$.  Every automorphism
therefore preserves $f$ up to a base Möbius transformation.  The two
multiplicities and the cusp distinguish its three special fibres, so
that transformation is the identity.

\item A holomorphic torus automorphism is affine: its lifted derivative is
lattice-periodic, descends to a holomorphic matrix-valued function on a
compact torus, and hence is constant.  Integrating gives an affine lift.
The locally constant integral linear part must commute with
\[
 A_1=\begin{pmatrix}1&0&0&0\\6&0&1&0\\-6&-1&-1&0\\-2&1&0&1\end{pmatrix},\quad
 A_2=\begin{pmatrix}1&0&0&0\\0&0&-1&0\\-6&1&0&0\\3&0&1&1\end{pmatrix}.
\]
Solving $PA_j=A_jP$ over $\mathbb Q$ gives
$P=aI+bE_{41}$; the complete rational linear elimination is in
\texttt{code/part5/centralizer.py}.  Integrality gives $a,b\in\mathbb Z$,
and invertibility gives $a=\pm1$.  The equation $H\Pi=\Pi P$ first gives
$H=aI$ from the identity columns and then $b=0$ from the first column.

\item The linear part $-I$ is excluded by the first multiple fibre.  Normalize
the root base change $t=s^3$.  The automorphism lifts, fixing $s$; on the
central covering torus it is $z\mapsto-z+b$.  Commutation with the affine
deck transformation gives
\[
 (I-A_1)b=2v_1/3\pmod\Lambda.
\]
The torus character $\gamma:\Lambda_{\mathbb R}/\Lambda\to\mathbb R/\mathbb Z$
annihilates $I-A_1$, whereas the right side has character $2l_1/3$.
Freeness implies $3\nmid l_1$, a contradiction.  Thus every automorphism
is fibrewise translation.

\item We compute all components of the translation group.  Let $j$ be the
inclusion of the cusp-punctured orbifold and set
$\mathcal H=j_*\Lambda$.  Its cusp stalk is
$\ker(M_0-I)=\mathbb Z\widehat w\oplus\mathbb Z\widehat\delta$.
Let $\mathcal A$ be the sheaf of fibrewise translations extending to
the local compactification.  There is an exact sequence
\begin{equation}\label{eq:P5exponential}
0\longrightarrow\mathcal H\longrightarrow V\longrightarrow\mathcal A
\longrightarrow0.
\end{equation}
For the local check at the cusp, use the toroidal covering whose
component lattice is $\Gamma=\Lambda/\Lambda_{\rm tor}$.  A translation
acts identically on the nearby fibre lattice and on $\Gamma$, and
therefore lifts.  Its meridian remains bounding a disc through a smooth
point of $W$, which verifies the same assertion for the full local
cover.  On the punctured torus the lift has form $x\mapsto a(q)x$.
It sends the component $E_0$ to $E_v$; the index-one map
$B_0:\Gamma\to\mathbb Z^2$ is an isomorphism, so composing with the
unique corresponding deck shear makes it preserve $E_0$.  It then
preserves the torus $E_0\setminus\partial E_0$.  Evaluation at $(1,1)$
in the chart $(\mathbb C^*)^2\times\Delta_q$ shows that $a_i(q)$ are
holomorphic units at zero, and hence have logarithms.  Conversely their
multiplication extends on every toric chart and commutes with the deck
shears.  The logarithm kernel consists precisely of those periods with
no logarithmic term, namely $\ker(M_0-I)$.  On the finite root charts
ordinary logarithms prove the same exact sequence with their finite
equivariance retained.

\item Here is the integral lattice calculation, including torsion.  Put
\[
 M_0=(A_1A_2)^{-1}=\begin{pmatrix}1&0&0&0\\0&1&0&0\\0&1&1&0\\-1&0&0&1\end{pmatrix}.
\]
A cocycle on the punctured orbifold is a pair $(x_1,x_2)$ with
$N_1x_1=N_2x_2=0$, where
\[
 N_1=\begin{pmatrix}3&0&0&0\\6&0&0&0\\-12&0&0&0\\0&2&1&3\end{pmatrix},\qquad
 N_2=\begin{pmatrix}4&0&0&0\\12&0&0&0\\-12&0&0&0\\0&2&2&4\end{pmatrix}.
\]
Extension over the cusp says that $x_1+A_1x_2$ vanishes in
$\operatorname{coker}(M_0-I)$, so its first two coordinates vanish.
The complete integral solution is
\[
 x_1=(0,-b,2b-3c,c),\qquad x_2=(0,-2a-b,b,a),\quad a,b,c\in\mathbb Z.
\]
The principal cocycle of $(g,u,w,d)$ has parameters
\[
 \begin{pmatrix}a\\b\\c\end{pmatrix}
 =\begin{pmatrix}3&0&1\\-6&1&-1\\-2&1&0\end{pmatrix}
  \begin{pmatrix}g\\u\\w\end{pmatrix}.
\]
This matrix has determinant $-1$; its inverse is explicitly
$g=c-a-b$, $u=3c-2a-2b$, $w=4a+3b-3c$.  Hence
$H^1(\widehat B,\mathcal H)=0$ integrally.  The independent total-complex
calculation in \texttt{code/part5/orbifold\_lattice\_complex.py} prints
every differential and checks the same vanishing.

\item Sequence \eqref{eq:P5exponential} now makes every global translation a
global logarithm in $H^0(\widehat B,V)=\mathbb C$.  Its kernel is
$\Lambda^{A_1,A_2}=\mathbb Z\widehat\delta$.  The flow extends by
commutation with the affine finite actions and by torus multiplication
on the cusp charts.  This proves existence, faithfulness, and
completeness of the asserted group.
\end{enumerate}
\end{proof}

\begin{theorem}\label{night:P5.6}\label{stmt:P5.6}
\StatementRecord{night:P5.6}{proved}{Historical events: GPT-6 Astra (OpenAI), ChatGPT Pro mode, effort Pro, S6\_06\_BLIND, 16--18 November 2026: targeted proof check. Also checked in this writing. Pre-continuation writing events: GPT-6 Astra (OpenAI), ChatGPT Pro mode, effort Pro, S6\_10\_PAPERS, 18 November 2026: targeted proof check. Continuation metadata events: GPT-6 Astra (OpenAI), ChatGPT Pro mode, effort Pro, S6\_10\_PAPERS continuation, 19 November 2026: used and rechecked in this run; status-vocabulary reconciliation; metadata reconciliation only; no new mathematical check. \emph{Pre-continuation recorded proof scope.} Read the full cp25 re-derivation and re-derived the local syzygy/root-cohomology, KS period equation, integral translation-sheaf, or order-two obstruction calculation stated in the respective proof. These are credited DGKX comparisons; the event does not assert an independent historical discovery.}{the definitions and exact calculations in this statement}

For fixed admissible twist data, $X_c\simeq X_{c'}$ precisely when
$c-c'\in2\mathbb Z$.  This fixes the integral period of $c$; the
construction still permits an additive choice of $\beta^{\rm part}$.
\end{theorem}
\begin{proof}
\begin{enumerate}[label=\arabic*.,leftmargin=*]
\item The intrinsic pencil argument and the centralizer computation give
$P=aI+bE_{41}$ with $a=\pm1$ and $b=a(c-c')$.  The same root-fibre
character excludes $a=-1$.  Thus $b=c-c'$ is an integer.

\item Compute $H^2(\widehat B,\mathcal H)$ with a pair of pants, two root discs,
and a cusp disc.  The pair of pants has cochains $\Lambda\to\Lambda^2$
with differential $((A_1-I),(A_2-I))$.  A root disc has the periodic
cyclic resolution with alternating differentials $A_j-I,N_j$; the
cusp disc has $\ker(M_0-I)$ in degree zero.  An annulus has
$\Lambda\xrightarrow{A_j-I}\Lambda$.  The degree-one restriction from
the pair of pants is
\[
 (x_1,x_2)\longmapsto(x_1,x_2,-M_0(x_1+A_1x_2)).
\]
The total-complex ranks in degrees zero through three are $14,28,20,8$.
The script \texttt{orbifold\_lattice\_complex.py} prints its three
integer differentials, checks their zero products, and computes
$H^0=\mathbb Z$, $H^1=0$, $H^2=\mathbb Z\oplus\mathbb Z/2$.

\item The degree-two presentation also reduces by hand.  The invariant
lattices have bases $(\widehat\delta,\varepsilon_j)$, where
$\varepsilon_1=(1,2,-4,0)$ and $\varepsilon_2=(1,3,-3,0)$.
The five generators $(\delta_1,\varepsilon_1,\delta_2,\varepsilon_2,h_c)$
have relations
\begin{equation}\label{eq:P5H2}
 \delta_1=0,\qquad 2\delta_2=0,\qquad
 3\varepsilon_1=h_c=4\varepsilon_2.
\end{equation}
For the constant sheaf $\mathbb Z$, the same computation has generators
$z_1,z_2,h_c$ and relations $3z_1=h_c=4z_2$; its primitive generator is
$z_1-z_2$.  Inclusion $\mathbb Z\widehat\delta\to\mathcal H$ sends it
to the nonzero order-two class $\delta_2$.  The character $\gamma$
sends $\varepsilon_j$ to $z_j$ and kills $\delta_j$.

\item Both $H^1(\widehat B,V)$ and $H^2(\widehat B,V)$ vanish, by
$\pi_*V=\mathcal O\oplus\mathcal O(-1)$ and the dimension of the coarse
curve.  Consequently \eqref{eq:P5exponential} identifies compactified
translation torsors $H^1(\widehat B,\mathcal A)$ with
$H^2(\widehat B,\mathcal H)$.  This identification takes logarithms of
local translation transitions and their lattice-valued Cech
coboundaries; the two vanishings also reverse that construction.

\item The cyclic norm of the translation $v_j/m_j$ is $v_j$, while the cusp
collar has the opposite boundary orientation.  Thus the torsor
character is
\[
 l_1z_1+l_2z_2-l_0h_c
 =(4l_1+3l_2-12l_0)(z_1-z_2)=-p(z_1-z_2).
\]
The shear $S_b-I=b\widehat\delta\otimes\gamma$ changes its class by
$-bp\delta_2$.  Freeness forces $l_2$ odd and thus $p$ odd.  The class
is preserved precisely for even $b$.

\item For an integral shear the Jacobian charts themselves identify: the
period lattices differ by $S_b$, and the bounded cusp matrix changes by
an integer in its $\beta$ entry.  Every toric unit
$\exp(2\pi iC\lambda)$ is therefore unchanged.  The fan, lattice, and
toric gluings are the same, and on root discs the basis change
identifies the smooth Jacobian and its cyclic action.  Equality of the
torsor classes now gives a holomorphic isomorphism for the fixed sheaf
$\mathcal A$, which consists precisely of translations extending over
the cusp.  It extends on all three special charts.  Conversely every
isomorphism preserves the class, so forces $b$ even.
\end{enumerate}
\end{proof}